\documentclass[10pt,a4paper]{article}
\usepackage[T1]{fontenc}
\usepackage{lmodern}
\usepackage{amsmath,amssymb,amsthm,mathtools}
\usepackage[margin=24mm]{geometry}
\usepackage{microtype}
\usepackage{booktabs,array,longtable}
\usepackage{xcolor}
\usepackage{hyperref}
\usepackage{fancyhdr}
\hypersetup{colorlinks=true,linkcolor=black,citecolor=black,urlcolor=blue,pdfauthor={Matthew J. Colbrook}}
\pagestyle{fancy}
\fancyhf{}
\fancyhead[L]{\small Partial result / MI-08}
\fancyhead[R]{\small 11 September 2026}
\fancyfoot[C]{\thepage}
\setlength{\headheight}{14pt}
\setlength{\parskip}{4pt}
\setlength{\parindent}{0pt}
\newtheorem{theorem}{Theorem}[section]
\newtheorem{proposition}[theorem]{Proposition}
\newtheorem{lemma}[theorem]{Lemma}
\newtheorem{corollary}[theorem]{Corollary}
\theoremstyle{remark}
\newtheorem{remark}[theorem]{Remark}
\newcommand{\M}{\mathbb M}
\newcommand{\C}{\mathbb C}
\newcommand{\R}{\mathbb R}
\newcommand{\tr}{\operatorname{tr}}
\newcommand{\diag}{\operatorname{diag}}
\newcommand{\spec}{\operatorname{spec}}
\newcommand{\rank}{\operatorname{rank}}
\newcommand{\abs}[1]{\lvert #1\rvert}
\newcommand{\norm}[1]{\lVert #1\rVert}
\newcommand{\opnorm}[1]{\lVert #1\rVert_{\infty}}
\newcommand{\schatten}[2]{\lVert #1\rVert_{#2}}
\newcommand{\symmod}{\mathcal S}
\newcommand{\maxmod}{\mathcal M}
\newcommand{\draftnotice}{\begin{center}\small\textit{Proposed mathematical resolution. Complete proof supplied; not submitted or peer reviewed.}\end{center}}

\usepackage{xurl}
\setlength{\emergencystretch}{3em}
\allowdisplaybreaks[1]

\title{MI-08: Partial Result}
\author{Matthew J. Colbrook\\
\small Department of Applied Mathematics and Theoretical Physics\\
\small University of Cambridge, Cambridge, United Kingdom\\
\small\href{mailto:m.colbrook@damtp.cam.ac.uk}{m.colbrook@damtp.cam.ac.uk}}
\date{11 September 2026}
\begin{document}
\maketitle
\textbf{Repository status: Partially resolved.}
The complete argument received an independent Codex-agent PASS review
for the stated partial result only.
The original draft was AI-assisted; the reviewed proof below is unchanged.
This is independent agent verification, not external human peer review or
formal proof certification. \href{https://github.com/ajt60gaibb/OpenProblemsInNLA/blob/main/references/colbrook-matrix-2026-09-11/verification/reviews/MI-08-review.md}{Detailed review and scope}.
\par\medskip
% BEGIN REVIEWED PROOF
\section{MI-08: fixed and adaptive orthogonal pinching have the same length}
\label{sec:mi08}
This section is partial progress on MI-08, not a determination of the length in every dimension. It resolves the comparison between the fixed and input-dependent versions in \cite[Questions 4.8--4.9]{BL2026} and determines additional exact dimensions.

Let $\Delta(X)=\diag(x_{11},\ldots,x_{dd})$ for real $d\times d$ matrices. Let $f(d)$ be the least $q$ for which there are fixed orthogonal matrices $O_1,\ldots,O_q$ satisfying
\[
 \Delta(X)=\frac1q\sum_{i=1}^qO_iXO_i^T\quad\hbox{for every real }X.
\]
Let $a(d)$ be the least $q$ for which the same identity holds for every real $X$, allowing the matrices $O_i$ to depend on $X$. Finally, let
\[
 h(d)=\min\{q:\ \exists H\in\{-1,1\}^{q\times d},\ H^TH=qI_d\}.
\]
Such a $q$ exists, for example by the Sylvester construction of sign matrices of power-of-two order.

\begin{theorem}\label{thm:mi08}
For every positive integer $d$, one has
\[
 f(d)=a(d)=h(d).
\]
In particular, $f(d)=a(d)=12$ for $9\leq d\leq12$.
\end{theorem}
\begin{proof}
If $H$ realizes $h(d)$, set $O_i=\diag(h_{i1},\ldots,h_{id})$. Orthogonality of the columns of $H$ gives
\[
 \frac1q\sum_iO_iXO_i^T=\Delta(X)
\]
entry by entry, for every real $X$. Thus $a(d)\leq f(d)\leq h(d)$.

For the reverse inequality, choose $D=\diag(1,2,\ldots,d)$ and the skew-symmetric matrix $K$ with $k_{ij}=1$ for $i<j$ and $k_{ij}=-1$ for $i>j$. Set $X=D+K$. Suppose $q$ orthogonal matrices, even chosen specifically for this $X$, satisfy the required identity. Taking symmetric parts yields
\[
 \frac1q\sum_iO_iDO_i^T=D.
\]
Every term has the same Frobenius norm as $D$. Therefore
\begin{align*}
 \sum_i\norm{O_iDO_i^T-D}_F^2
 &=2q\norm{D}_F^2-2\tr\left(D\sum_iO_iDO_i^T\right)=0.
\end{align*}
Each $O_i$ consequently commutes with $D$. Since the diagonal entries of $D$ are distinct, $O_i$ is a diagonal sign matrix. Taking skew-symmetric parts now shows that, for every $j\ne l$,
\[
 0=\sum_i (O_i)_{jj}(O_i)_{ll}k_{jl}.
\]
All $k_{jl}$ are nonzero. The $q\times d$ matrix of these signs therefore has orthogonal columns, so $q\geq h(d)$. This proves $a(d)\geq h(d)$ and establishes the equality.

For $d\geq3$, any admissible $q$ is divisible by four. Indeed, multiply rows by signs to make the first column all ones. The second and third columns are both balanced and orthogonal; the four possible sign pairs must each occur $q/4$ times. Also $q\geq d$ by rank. Thus $h(d)\geq12$ for $9\leq d\leq12$.

For completeness, an order-twelve sign matrix can be constructed entirely explicitly. In the field $\mathbb F_{11}$, put $\chi(0)=0$, let $\chi(x)=1$ for nonzero squares, and let $\chi(x)=-1$ otherwise. Let $Q_{ij}=\chi(i-j)$ for $i,j\in\mathbb F_{11}$. Then $Q^T=-Q$, $Q\mathbf1=0$, and $QQ^T=11I-J$. The last identity can also be checked directly using the square residues $\{1,3,4,5,9\}$; the accompanying verifier does so with integer arithmetic. Hence
\[
 H=\begin{pmatrix}1&\mathbf1^T\\\mathbf1&Q-I\end{pmatrix}
 \quad\hbox{satisfies}\quad HH^T=12I_{12}.
\]
Taking any $d$ columns gives $h(d)\leq12$ for $d\leq12$, proving the claimed values.
\end{proof}

\begin{remark}[What remains]
The general value of $h(d)$ is not determined here. In particular, $h(4m)=4m$ is equivalent to the existence of a Hadamard matrix of order $4m$. The value $h(9)=12$ also rules out a universal dichotomy between $d$ and the next power of two. The adaptive comparison above uses all real matrices, including nonsymmetric ones, exactly as in the stated question; it is not a claim about a restriction to symmetric inputs only.
\end{remark}

% END REVIEWED PROOF
\begin{thebibliography}{1}

\bibitem{BL2026}
Jean-Christophe Bourin and Eun-Young Lee.
\newblock Averages over matrix unitary orbits and spectral order.
\newblock arXiv:2606.15624, 2026.
\newblock Version consulted: v2; Question 3.11, Conjecture 3.12, and Questions
  4.8--4.9.
\newblock URL: \url{https://arxiv.org/abs/2606.15624}.

\end{thebibliography}

\end{document}
